% Generated draft: requires mathematical and source review.
\documentclass[11pt]{article}
\usepackage[margin=25mm]{geometry}
\usepackage{fontspec}
\setmainfont{Latin Modern Roman}
\usepackage{amsmath,amssymb,mathtools,unicode-math}
\setmathfont{Latin Modern Math}
\usepackage{enumitem,needspace,xurl,hyperref}
\setlength{\parindent}{0pt}
\setlength{\parskip}{5pt}
\setlength{\emergencystretch}{3em}
\newcommand{\N}{\mathbb{N}}
\newcommand{\Z}{\mathbb{Z}}
\newcommand{\Q}{\mathbb{Q}}
\newcommand{\R}{\mathbb{R}}
\newcommand{\C}{\mathbb{C}}
\newcommand{\F}{\mathbb{F}}
\newcommand{\A}{\mathbb{A}}
\newcommand{\PP}{\mathbb{P}}
\newcommand{\E}{\mathbb{E}}
\newcommand{\eps}{\varepsilon}
\newcommand{\dd}{\,\mathrm{d}}
\begin{document}
\section*{UnsolvedMath ID: 800122}
\textbf{Source problem code:} catalogue entry
\subsection{Definitions and mathematical statement}
The supplied question uses $B\bar\Gamma_q$ without defining it or specifying differentiability, transverse orientation, or the range of codimensions $q$. The following formulation is conditional on the conventional smooth, unoriented interpretation of this notation.

For a nonnegative integer $q$, let $\Gamma_q$ denote the Haefliger topological groupoid of germs of local smooth diffeomorphisms of $\R^q$, and let $B\Gamma_q$ be its classifying space. The normal vector bundle of the universal Haefliger structure determines a map
\[
\nu_q:B\Gamma_q\longrightarrow BO(q),
\]
where $BO(q)$ classifies real vector bundles of rank $q$. In this interpretation, define
\[
B\bar\Gamma_q:=\operatorname{hofib}_{*}(\nu_q),
\]
with the homotopy fiber taken over the point classifying the trivial bundle.

The question, identified in the supplied source as the \emph{Haefliger--Thurston conjecture}, asks whether, for every $q$ in the intended codimension range, $B\bar\Gamma_q$ is $2q$-connected. Explicitly, does one have
\[
B\bar\Gamma_q\text{ is path-connected},\qquad
\pi_i(B\bar\Gamma_q,x)=0
\quad\text{for every }x\in B\bar\Gamma_q
\text{ and every integer }1\leq i\leq 2q?
\]
The smooth, unoriented interpretation and the intended range of $q$ require confirmation from the source; they are not specified by the supplied excerpt.

The source places this question in a discussion connecting homotopical triviality of Haefliger classifying spaces with existence of foliations and homological acyclicity of diffeomorphism groups. It describes negative answers as potentially revealing obstructions to Reeb's fundamental problem beyond Bott's obstructions and secondary characteristic classes such as the Godbillon--Vey invariant. These are contextual motivations supplied by the source, without precise hypotheses or implication statements in this record.
\subsection{Short English statement}
Under the conventional smooth, unoriented interpretation, the question asks whether the homotopy fiber of the universal normal bundle map $B\Gamma_q\to BO(q)$ is path-connected and has vanishing homotopy groups through degree $2q$, for every intended codimension $q$. The excerpt leaves the regularity convention and codimension range unspecified. Its label \emph{open} is supplied registry metadata, not an independently verified status assessment.
\subsection{Sources}
\begin{itemize}
\item Local UnsolvedMath registry, permanent ID $800122$, exact source problem code \texttt{catalogue entry}, upstream revision \texttt{d371ee917ac9647fbb06cd6fc529deec6f3dda82}. The record attributes the source list to Nariman and Yazdi, gives the year $2025$, and labels the problem \emph{open}.
\item Nariman and Yazdi, \emph{Problem List on Foliations and Diffeomorphism Groups} (2025), Question 63: \url{https://arxiv.org/abs/2503.18274}. The supplied attribution/context is ``Haefliger--Thurston conjecture.'' The record reports access on 2026-10-16 and extraction by \texttt{arxiv\_latex}; these provenance details are supplied and have not been independently verified here.
\item The supplied excerpt mentions Reeb's fundamental problem, Bott's obstructions, and secondary characteristic classes, including the Godbillon--Vey invariant, as motivation. No separate references for these contextual claims or for the notation $B\bar\Gamma_q$ are supplied.
\end{itemize}
\end{document}
